\newpage
\section{Perturbations} 

\subsection{List of Datasets} \label{app:perturbation_datasets}

\paragraph{Passages}
\begin{itemize}
    \item \hfbutton{https://huggingface.co/datasets/allegrolab/passages_gutenberg_popular} \textbf{Gutenberg Popular} are passages sampled from the popular books from the Gutenberg corpus \citep{gerlach2018standardizedprojectgutenbergcorpus}. Due to studies like \cite{kirchenbauer2024lmd} which show pretraining data density affects memorization, we stratify two Gutenberg splits based on download counts. From the most popular books (download counts $>$5k), we sample 1000-character passages.
    
    \item \hfbutton{https://huggingface.co/datasets/allegrolab/passages_gutenberg_unpopular} \textbf{Gutenberg Unpopular} are sampled passages from the unpopular books from the Gutenberg corpus \citep{gerlach2018standardizedprojectgutenbergcorpus}. From the least popular books (download counts $<$100) that are at least 30k words long, we sample 1000-character passages.
    
    \item \hfbutton{https://huggingface.co/datasets/allegrolab/passages_wikipedia} \textbf{Wikipedia} are passages sampled from our crawl of Wikipedia articles. We begin our crawl at the Wikipedia pages "2023" and "2024". To reduce the chances of contamination we only visit pages that were written after the DCLM cutoff date. After filtering out articles without text (e.g. lists), we end up with 1500 articles. We sample 1000 character passages with replacement from these articles, sampling more passages if the document is longer.
\end{itemize}

\paragraph{Paraphrases}
\begin{itemize}
    \item \hfbutton{https://huggingface.co/datasets/allegrolab/paraphrases_mrpc} \textbf{MRPC} \citep{dolan2005automatically} are paraphrases where the source sentences are drawn from news articles headlines. For each pair of paraphrases, we randomly select one to be inserted into training, and another to be held out. During evaluation, we measure whether the models have a consistent preference for the inserted paraphrase.
    
    \item \hfbutton{https://huggingface.co/datasets/allegrolab/paraphrases_paws} \textbf{PAWS} \citep{zhang-etal-2019-paws} is a dataset of paraphrases generated by rule-based word swaps and backtranslation. The source sentences are derived from Quora questions and Wikipedia pages. Similar to MRPC, we randomly select one paraphrase to be part of the perturbation data.
\end{itemize}

\paragraph{Biographies}
\begin{itemize}
    \item \hfbutton{https://huggingface.co/datasets/allegrolab/biographies_yago} \textbf{YAGO}: We synthetically generate biographies of fictional people using distributions computed from YAGO, a real-world knowledge graph \citep{tanon_yago_2020}. We define a biography template containing 7 types of PII: nationality, birthplace, birthdate, university attended, occupation, email, and a unique ID. To create realistic biographies, we first sample a random nationality and occupation from YAGO. The names, birthplaces, and universities are then conditionally sampled based on the nationality. Finally, birthdates, emails, and UUIDs are randomly sampled. Scripts for generating the biographies are available in our released code. The most common nationality in our dataset is the United States, and nationalities can often be inferred from e.g. the birthplace, as they are correlated information.
    
    \item \hfbutton{https://huggingface.co/datasets/allegrolab/biographies_ecthr} \textbf{ECtHR} \cite{pilan2022text} introduces a text anonymization benchmark based on a collection of European court records annotated for personally identifiable information. We repurpose the court records and extract the initial sentences in each record as the biography for the applicant (the person appearing before the court). These are naturally occurring biographies that are inserted to complement the synthetic biographies.
\end{itemize}

\paragraph{Chats}

\begin{itemize}
    \item \hfbutton{https://huggingface.co/datasets/allegrolab/chats_personachat} \textbf{Personachat} \citep{zhang-etal-2018-personalizing} is a dataset where two crowdworkers engaged in a conversation based on the personas assigned to them. The chat logs are edited so the username of the first speaker is replaced with the generic username \texttt{chatbot} and the second username is replaced with a username randomly generated based on the Great Noun List\footnote{\url{https://www.desiquintans.com/nounlist}}. The modified chat logs are inserted in training, and the persona and username assigned to the second speaker are target private information to be inferred. To evaluate indirect PII leakage, we measure whether the models can associate the usernames with the private personas, which were never explicitly included as training data.
\end{itemize}

\paragraph{Standard test sets}
\begin{itemize}
    \item \hfbutton{https://huggingface.co/datasets/allegrolab/testset_popqa} \textbf{PopQA} \citep{mallen-etal-2023-trust} is an open-ended question answering dataset that evaluates the world knowledge of a model. To contaminate the task, we insert questions followed by the answer. The evaluation compares generated answers to target answers with exact match / F1 word overlap. 
    
    \item \hfbutton{https://huggingface.co/datasets/allegrolab/testset_winogrande-infill} \textbf{Winogrande-Infill} \citep{sakaguchi_2021_winogrande} is a binary  pronoun resolution task where the model is given a context and asked to determine which entity a pronoun refers to. Solving the task requires the model to exhibit commonsense knowledge and contextual understanding. Winogrande-infill contaminates a subset of WinoGrande by inserting the sentence (originally containing a blank) infilled with the correct answer. Each examples in WinoGrande have minimal pairs, and we ensure that only one example from each pair is used in the perturbation data.

    \item \hfbutton{https://huggingface.co/datasets/allegrolab/testset_winogrande-mcq} \textbf{Winogrande-MCQ} is a second contamination variant for Winogrande. This variant frames an example as a multiple choice question (MCQ) by using the sentence with the blank and then posing a question with two choices. We insert the question followed by the correct answer in the corpus. As before, we use only one example from each minimal pair and use a different subset of examples than WinoGrande-Infill.

    \item \hfbutton{https://huggingface.co/datasets/allegrolab/testset_MMLU} \textbf{MMLU} \citep{hendrycks2021measuring} is a 4-way multiple choice question answering dataset that covers 57 different domains and tasks, evaluating both world knowledge and problem-solving capabilities. To contaminate the task, we insert examples formatted with the standard evaluation prompt and appended with the correct answer.
    
    \item \hfbutton{https://huggingface.co/datasets/allegrolab/testset_hellaswag} \textbf{HellaSwag} \citep{zellers-etal-2019-hellaswag} is a 4-way multiple choice commonsense reasoning dataset, where the model is required to understand implicit context and common knowledge in order to correctly select the continuation to a context. Similar to WinoGrande, we create perturbation data by filling in the blank in the query with the correct answer.
    
    \item \hfbutton{https://huggingface.co/datasets/allegrolab/testset_piqa} \textbf{PIQA} \citep{Bisk_Zellers_Lebras_Gao_Choi_2020} is a binary multiple choice question answering dataset that requires the model to use physical commonsense reasoning to answer correctly. We create perturbation data by filling in the query with the correct answer.
\end{itemize}

\paragraph{New test sets}

\begin{itemize}
    \item \hfbutton{https://huggingface.co/datasets/allegrolab/testset_ellie} \textbf{ELLie} \citep{testa-etal-2023-understand} tests the language model's understanding of ellipsis. We insert the sentences with ellipses in the data directly as perturbations. For evaluation, we use the GPT prompt format defined for each example.
    \item \hfbutton{https://huggingface.co/datasets/allegrolab/testset_munch} \textbf{MUNCH} \citep{tong-etal-2024-metaphor} tests a language model's ability to differentiate between apt and inapt usage of metaphors in a sentence. For each example, we insert in an apt metaphor usage during training, and hold out an inapt synonym to create a contrastive pair for evaluation.
\end{itemize}

\subsection{Perturbation Statistics}
\label{app:perturbation-data-statistics}

\begin{table}[h]
    \centering
    \caption{\textbf{Percentage of training data overwritten by duplicated perturbation data.} These calculations depend on the selected sequence length of 2048 tokens and training batch size of 1024 sequences.}
    \label{tab:perturbation-data-stats}
    \begin{tabular}{lcc}
        \toprule
        & \textbf{100B} & \textbf{500B} \\
        \midrule
        \textbf{Tokens Modified} & 0.08\% & 0.016\% \\
        \textbf{Sequences Modified} & 1.67\% & 0.34\% \\
        \textbf{Avg. Perturbations per Batch} & 17 & 3.4 \\
        \bottomrule
    \end{tabular}
\end{table}

For each perturbation type, we sought to (1) insert different levels of duplications to induce a range of memorization and (2) duplicate enough examples at each level to achieve precise memorization estimates for that level. Based on initial experiment of 1B models, we find the range of duplications $\{0,1,4,16,64,256\}$ to induce a range of memorization. For smaller datasets, we only duplicate powers of 16, up to 256. 

For the 0 and 1 duplicate levels, we aimed to insert more than 1000 examples (derived from a binomial power calculator), which yields small error bars. At the highest duplication level (256), we typically insert only 1/10th of examples at the lowest duplication level (1). When an example is highly duplicated and strongly memorized, there is typically low entropy in the model predictions so the resulting error bars over less examples are still small. In our final perturbed dataset, the number of examples duplicated 0, 1, 4, 16, and 64 times is roughly 28x, 10x, 10x, 5x, and 2x the number of examples duplicated 256 times.

% give a table of the duplicate levels of each dataset and the number of examples


% For 100B
% > Inserting 818393 samples.
% > Perturbing 1.6650248209635417% of the training samples.
% > Perturbing 1704.9854166666667% of the training batches (expectation).
% For 500B
% > Inserting 818393 samples.
% > Perturbing 0.3350993350366876% of the training samples.
% > Perturbing 343.1417190775681% of the training batches (expectation).

% Token Count
% Testset - 25.5M
% Copyright - 36.8M
% Privacy - 17.6M

\subsection{Decontamination}
\label{app:decontamination}

To ensure accurate duplication counts for our perturbations, we decontaminate the documents and perturbation data in two phases, depending on the length of the perturbations. For perturbations \textit{longer than 10 tokens}, we decontaminate the training data. We build an Infini-gram index \citep{liu2024infinigram}, enabling fast queries for exact matches over all training documents. Here, we query and remove training documents that have large n-gram overlaps with our perturbations \citep[similar to][]{brown_2020_language}. The threshold is chosen conservatively to avoid spurious matches and identify duplicated test sets. For documents up to 40 tokens, we check for exact matches with the full document. For documents longer than 40 tokens ($n>40$), we search for matches using $n/2$-grams with a stride of $n/4$ tokens. For \textit{test set perturbations} (usually very short), removing matching training documents risks discarding too many documents. Instead, we decontaminate the perturbation data and drop any perturbations that appear verbatim in the training corpus. When applicable, we use multiple query formats to identify matches. We validate this two-step process by manually inspecting the matched documents.
